
\documentclass[../../../physics-standard_answers_all.tex]{subfiles}

\begin{document}

\setlength{\abovedisplayskip}{2mm}
\setlength{\belowdisplayskip}{0mm}

\begin{enumerate}

    \item 運ぶ際の電位の変化を$\Delta \phi$とすれば，陽子の位置エネルギー変化
    $e\Delta \phi$に等しい仕事が必要である．
    %
    \begin{align*}
        & W_{\mathrm{PF}} = 1.6 \times 10^{-19} \, \mathrm{C} \times (14-6) \, \mathrm{V}
        \fallingdotseq \uwave{1.3 \times 10^{-18} \, \mathrm{J}} \;,
        \\
        & W_{\mathrm{PG}} = 1.6 \times 10^{-19} \, \mathrm{C} \times (6-6) \, \mathrm{V}
        = \uwave{0 \, \mathrm{J}} \;,
        \\
        & W_{\mathrm{PH}} = 1.6 \times 10^{-19} \, \mathrm{C} \times (4-6) \, \mathrm{V}
        = \uwave{-3.2 \times 10^{-19} \, \mathrm{J}} \;.
    \end{align*}

    \item 運ぶ際の電位の変化を$\Delta \phi$とすれば，電子の位置エネルギー変化
    $(-e) \Delta \phi$に等しい仕事が必要である．
    %
    \begin{align*}
        & 
        \begin{aligned}
        W^*_{\mathrm{PF}} &= -1.6 \times 10^{-19} \, \mathrm{C} \times (14-6) \, \mathrm{V} \\
         & \fallingdotseq \uwave{-1.3 \times 10^{-18} \, \mathrm{J}} \;,
        \end{aligned}
        \\
        & W^*_{\mathrm{PG}} = -1.6 \times 10^{-19} \, \mathrm{C} \times (6-6) \, \mathrm{V}
        = \uwave{0 \, \mathrm{J}} \;,
        \\
        & W^*_{\mathrm{PH}} = -1.6 \times 10^{-19} \, \mathrm{C} \times (4-6) \, \mathrm{V}
        = \uwave{3.2 \times 10^{-19} \, \mathrm{J}} \;.
    \end{align*}

    \item 電気力線は等電位線に直交し，高電位から低電位へ向かう向きである．        
        %
        \par \vspace{0.2\baselineskip}
        {\centering
            \includegraphics
            {\subfix{fig/fig_em_ef_04/fig_em_ef_04_03.pdf}}
        \par}
        \vspace{0.2\baselineskip}

    \item 電場は電位の空間勾配に等しいので，
        %
        \begin{align*}
            E_\mathrm{P} \fallingdotseq \frac{2 \, \mathrm{V}}{0.53 \, \mathrm{m}}
            = 3.77  \cdots \, \mathrm{V/m}
            \fallingdotseq \uwave{3.8 \, \mathrm{V/m}} \;.
        \end{align*}

\end{enumerate}

\end{document}
